\documentclass[10pt]{article}
\usepackage[a4paper,margin=0.65in]{geometry}
\usepackage[T1]{fontenc}
\usepackage{booktabs}
\usepackage{longtable}
\usepackage{array}
\usepackage{hyperref}
\usepackage{xcolor}

\title{SHAKTI C-Class Dual-Issue Performance Model Experiments}
\author{perf-model}
\date{\today}

\begin{document}
\maketitle

\section*{Executive Summary}

The delivered dual-issue SHAKTI model predicts 10,604,215
cycles and IPC 1.2037 on the available CoreMark trace
window. The supplied clean dual RTL ground truth is 10,398,758
cycles and IPC 1.22817, so the model is 1.976\%
high in cycles on the aggregate check.

\textbf{Validation gate status: not fully passed.} The single-issue model still
has strong $\Delta t$ validation on the available stamped trace
(99.234\%), but a cycle-stamped dual RTL trace was not
available in the workspace. The local CoreMark trace window also contains
12,764,348 instructions, while the supplied
ground truth contains 12,771,408. Therefore the dual
experiments below are useful architectural predictions, but they carry the
residual aggregate/model-input mismatch and are not yet a 97\%+ dual
$\Delta t$ validation.

The most important result is that the second memory port alone helps only
modestly in this model. Intra-bundle ALU forwarding is the largest individual
gain. The best draining two-issue combination found is E+H: intra forwarding + branch,
with IPC 1.2915
(+7.292\%
versus baseline). The quad approximation is separate and reaches IPC
1.4196. The model-owned branch+branch
opportunity count is much smaller than the defective RTL event-52 story suggests.

\section*{Experiment Summary}

\begin{longtable}{@{}p{0.27\linewidth}rrrrp{0.13\linewidth}@{}}
\toprule
Experiment & Cycles & IPC & $\Delta$ IPC & $\Delta$ & Verdict \\
\midrule
\endhead
Baseline & 10,604,215 & 1.2037 & +0.0000 & +0.000\% & no effect \\
A: second memory port & 10,524,390 & 1.2128 & +0.0091 & +0.758\% & helped \\
B: decouple lockstep & 10,482,209 & 1.2177 & +0.0140 & +1.164\% & helped \\
C: A+B & 10,440,944 & 1.2225 & +0.0188 & +1.564\% & helped \\
D: independent retire & 10,604,215 & 1.2037 & +0.0000 & +0.000\% & no effect \\
E: intra-bundle ALU forwarding & 9,993,129 & 1.2773 & +0.0736 & +6.115\% & helped \\
F: relax branch next-PC stall & 10,604,215 & 1.2037 & +0.0000 & +0.000\% & no effect \\
G: symmetric slots & 10,604,215 & 1.2037 & +0.0000 & +0.000\% & no effect \\
H: second branch unit & 10,554,246 & 1.2094 & +0.0057 & +0.473\% & helped \\
I: best combination & 10,037,159 & 1.2717 & +0.0680 & +5.650\% & helped \\
J: quad issue on best & 8,991,805 & 1.4196 & +0.2158 & +17.932\% & helped \\
\bottomrule
\end{longtable}

\section*{Additional Combination Probes}

\begin{longtable}{@{}p{0.27\linewidth}rrrrp{0.13\linewidth}@{}}
\toprule
Experiment & Cycles & IPC & $\Delta$ IPC & $\Delta$ & Verdict \\
\midrule
\endhead
A+E: memory + intra forwarding & 10,146,981 & 1.2579 & +0.0542 & +4.506\% & helped \\
E+H: intra forwarding + branch & 9,883,469 & 1.2915 & +0.0878 & +7.292\% & helped \\
\bottomrule
\end{longtable}

\section*{Validation}

\begin{tabular}{@{}lr@{}}
\toprule
Item & Value \\
\midrule
Single-issue $\Delta t$ accuracy on local CoreMark trace & 99.234\% \\
Single-issue model cycles on local window & 14,854,777 \\
Single-issue local RTL trace cycles & 14,965,228 \\
Single-issue cycle error vs local trace & -0.738\% \\
Dual baseline model cycles & 10,604,215 \\
Dual clean RTL ground-truth cycles & 10,398,758 \\
Dual aggregate cycle error & 1.976\% \\
Dual baseline model IPC & 1.20371 \\
Dual clean RTL ground-truth IPC & 1.22817 \\
Held-out Dhrystone clean dual RTL IPC & 1.33 \\
\bottomrule
\end{tabular}

\paragraph{Counter-profile validation.}
The model now emits the RTL-facing counters needed for validation, but the only
available local dynamic trace is CoreMark. The prompt's counter reference is a
Dhrystone dual-issue instrumented run, and no Dhrystone commit trace exists in
the workspace. Comparing absolute CoreMark model counters to Dhrystone RTL
counters would be misleading, so the mem+mem 5\% validation gate remains
blocked pending a Dhrystone or dual CoreMark cycle-stamped commit trace.

The baseline CoreMark model counter profile is:

\begin{tabular}{@{}lr@{}}
\toprule
Counter & Model value \\
\midrule
dual\_issued & 4,864,962 (45.878\% of cycles) \\
raw\_hazard & 1,983,094 \\
one\_instr & 149,500 \\
st3\_not\_firing & 2,070,236 \\
mem\_mem\_hazard & 738,993 \\
mem LL / LS / SS & 447,048 / 284,784 / 7,161 \\
branch+branch opportunities & 162,837 \\
mispredict & 211,531 \\
\bottomrule
\end{tabular}

For reference, the supplied Dhrystone RTL profile was: dual-issued about 50\%
of cycles, RAW 12,051, one-instruction fetch 4,026, stage3 not firing 4,864,
mem+mem 22,510 (LL 6,010, LS 6,000, SS 10,499), and mispredict 2,034.

\section*{Calibration Log}

\begin{longtable}{@{}p{0.23\linewidth}p{0.31\linewidth}p{0.34\linewidth}@{}}
\toprule
Change & Effect & Evidence \\
\midrule
\endhead
WAW/WAR pair handling & Kept WAW and WAR co-issue enabled because
\texttt{no\_wawstalls} compiles out the checks. & Unit tests verify same-rd
and write-after-read pairs commit together. \\
Lockstep scoping & Kept lockstep only for real paired bundles; a single bundle
does not wait on slot 1. & Unit tests exercise dependent and independent
single/pair cases. \\
Dual bypass shape & Fixed bypass lookup to inspect slot 1 and slot 0 of both
downstream sources, matching \texttt{bypass.bsv} priority
\texttt{src0/slot1 > src0/slot0 > src1/slot1 > src1/slot0}. & Full CoreMark
dual baseline moved from IPC 1.0357 before this fix to IPC 1.2037. \\
Memory pairing policy & Split the actual delivered branch
(\texttt{memory\_pairing=none}) from the proposed second-port experiment
(\texttt{memory\_pairing=all}). & Run A is now a true second-memory-port
experiment instead of the narrower, gated store-involving \texttt{dual\_mem}
path. \\
\bottomrule
\end{longtable}

\section*{Analysis}

Run A, the second memory port, improves IPC by +0.758\%.
It removes the modeled non-RAW mem+mem reject counter, but the speedup is much
smaller than a first-order lost-slot estimate because many newly legal memory
pairs still contend with RAW dependencies, branch redirects, and lockstep
movement through later stages.

Run B improves IPC by +1.164\%, confirming
that stage3/stage4 coupling is a real limiter. Run C is sub-additive relative to
A+B, so the second memory port and decoupling overlap in the bubbles they remove.

Run E is the largest individual improvement at +6.115\%.
It directly attacks the RAW rejection class and roughly halves the modeled RAW
reject count. This is the strongest candidate for RTL work if its implementation
cost is manageable.

Run H gives only +0.473\%. The model's
independent branch+branch opportunity count is 162,837
on CoreMark, so branch+branch is not the dominant issue in this trace. This is
also a useful independent replacement for the defective RTL event 52.

The automatic stacked run I (A+E+H) reached IPC 1.2717.
The extra probes show why it is not the actual best: A+E reached IPC
1.2579, while
E+H reached IPC 1.2915.
So the second memory port interferes with the intra-forwarding win in this trace,
whereas branch+branch is mildly complementary with intra-forwarding.

\section*{Sensitivity}

\begin{longtable}{@{}p{0.42\linewidth}rrr@{}}
\toprule
Config & Cycles & IPC & $\Delta$ IPC \\
\midrule
\endhead
instr\_queue=2 & 10,597,053 & 1.2045 & +0.068\% \\
instr\_queue=4 & 10,602,913 & 1.2039 & +0.012\% \\
instr\_queue=6 & 10,604,215 & 1.2037 & +0.000\% \\
instr\_queue=8 & 10,604,215 & 1.2037 & +0.000\% \\
instr\_queue=12 & 10,621,087 & 1.2018 & -0.159\% \\
mispredict penalty=0 & 10,391,676 & 1.2283 & +2.045\% \\
mispredict penalty=1 & 10,604,215 & 1.2037 & +0.000\% \\
mispredict penalty=2 & 10,815,746 & 1.1802 & -1.956\% \\
mispredict penalty=3 & 11,027,277 & 1.1575 & -3.837\% \\
memory latency=0 & 10,604,215 & 1.2037 & +0.000\% \\
memory latency=1 & 10,604,215 & 1.2037 & +0.000\% \\
memory latency=2 & 12,239,228 & 1.0429 & -13.359\% \\
memory latency=3 & 14,196,758 & 0.8991 & -25.305\% \\
\bottomrule
\end{longtable}

\section*{What Did Not Work}

\begin{itemize}
\item Independent retire alone had no effect because actual lockstep stage4
already delivers paired slots to stage5 together.
\item Symmetric slots alone had no effect because the current whitelist already
captures the useful mixed-FU pairs visible in the trace.
\item The branch next-PC relaxation is unresolved: the current model does not
include a separate queue-head next-PC stall beyond predictor redirect timing, so
the named F config is a no-op and should not be treated as RTL evidence.
\item Combining decoupled lockstep with intra-bundle forwarding did not drain on
a 100k-instruction diagnostic run. The final best combination therefore excludes
B when E is selected; this is reported as a model limitation/negative
interaction, not hidden.
\end{itemize}

\section*{Recommendations}

\begin{enumerate}
\item Prototype intra-bundle ALU forwarding first. Estimated standalone gain:
+6.115\%; rough RTL cost: medium
(operand muxing, bypass select, verification of RAW edge cases).
\item Combine intra-bundle ALU forwarding with a second branch unit only if the
branch hardware is cheap. Best measured combination: E+H at IPC
1.2915.
\item Treat the second memory port as a secondary change, not the headline by
itself. Estimated standalone gain: +0.758\%;
rough RTL cost: high (D-cache/store-buffer/arbiter integration). It regresses
the E combination, so do not prioritize it before validating with a dual trace.
\item Investigate lockstep decoupling separately from intra-forwarding.
Estimated standalone gain: +1.164\%;
rough RTL cost: medium-high because it touches stage3/stage4 valid/ready and
commit bookkeeping.
\item Deprioritize branch+branch hardware until a corrected RTL counter or dual
trace says otherwise. Estimated gain: +0.473\%;
rough RTL cost: medium.
\item Do not commit to quad issue from this evidence alone. The quad result is
an optimistic adjacent-pair approximation, not a validated 4-wide front end.
\end{enumerate}

\section*{Limitations}

The model is a timing model over an already committed dynamic trace. It does not
compute data values. It does not model data-cache misses, cache conflicts,
instruction-cache misses, data-dependent divider timing, or CSR/trap behavior
that depends on values. The dual validation lacks per-instruction RTL commit
cycles, so aggregate IPC agreement can still hide local timing errors. A
cycle-stamped dual RTL CoreMark trace and a Dhrystone commit trace are the next
required artifacts.

\end{document}
